% !TEX root = ../main.tex
%% Sujet type bac n° 4 — sujet inédit OptiTAB (valeurs vérifiées : verifications/sujets_bac.py)

\sujetBac{4}{{L'eau oxygénée : oxydoréduction, cinétique, gaz parfait}/9,{Une lunette astronomique : optique, diffraction}/6,{Le château d'eau : statique et dynamique des fluides}/5}

\begin{exoBac}{9}{L'eau oxygénée}
L'eau oxygénée, solution aqueuse de peroxyde d'hydrogène $\ce{H2O2}$, est un antiseptique. Versée sur une plaie, elle
mousse : une enzyme du sang, la catalase, accélère fortement sa décomposition en eau et en dioxygène.
\begin{donnees}
Couples oxydant-réducteur : $\ce{O2}/\ce{H2O2}$ et $\ce{H2O2}/\ce{H2O}$. Volume molaire d'un gaz à \qty{0}{\degreeCelsius} sous \qty{1,013}{bar} :
\qty{22,4}{L.mol^{-1}}. $R=\qty{8,314}{J.mol^{-1}.K^{-1}}$ ; $T(\unit{K})=\theta(\unit{\degreeCelsius})+273$.
\end{donnees}

\partieExo{1}{Le titre d'une eau oxygénée}
\begin{docu}{1}{Une eau oxygénée « à 10 volumes »}
Un litre d'eau oxygénée à 10 volumes peut libérer, en se décomposant totalement, \qty{10}{L} de dioxygène, volume
mesuré à \qty{0}{\degreeCelsius} sous \qty{1,013}{bar}.
\end{docu}
\begin{questions}
\item Écrire les deux demi-équations électroniques, puis montrer que l'équation de la décomposition s'écrit
$\ce{2H2O2 -> 2H2O + O2}$. Quelle particularité présente le peroxyde d'hydrogène dans cette réaction ?
\item Montrer que la concentration en peroxyde d'hydrogène d'une eau oxygénée à 10 volumes vaut \qty{0,89}{mol.L^{-1}}.
\end{questions}

\partieExo{2}{Suivi cinétique de la décomposition}
\begin{docu}{2}{Le protocole}
Dans un flacon fermé, relié à un capteur de pression, on place $V=\qty{20,0}{mL}$ d'une solution de peroxyde d'hydrogène de
concentration $c_0=\qty{0,20}{mol.L^{-1}}$ ; le volume de gaz au-dessus de la solution vaut $V_\mathrm{g}=\qty{250}{mL}$ et la température
reste égale à \qty{25}{\degreeCelsius}. À $t=0$, on ajoute quelques gouttes d'une solution d'ions fer(III), qui catalysent la décomposition.
Le capteur mesure la surpression $\Delta P$ due au dioxygène formé.
\end{docu}
\schema{bac-s4-pression}
\begin{questions}[resume]
\item Établir le tableau d'avancement de la transformation. Montrer que l'avancement vaut $x=\dfrac{\Delta P\,V_\mathrm{g}}{R\,T}$.
\item Calculer l'avancement maximal et la surpression finale attendue. Comparer à la courbe.
\item Définir le temps de demi-réaction et déterminer sa valeur.
\item Les ions fer(III) sont-ils un catalyseur homogène ou hétérogène ? Pourquoi la vitesse de formation du dioxygène
diminue-t-elle au cours du temps ?
\item \capacite{Capacité numérique}\initiative Le programme suivant teste l'hypothèse d'une loi de vitesse d'ordre 1, pour laquelle
$[\ce{H2O2}]=c_0\,\mathrm{e}^{-kt}$, soit $\Delta P=\Delta P_\mathrm{max}\left(1-\mathrm{e}^{-kt}\right)$. Interpréter les valeurs affichées et conclure.
\begin{codecote}[0.66]
\begin{pythoncol}
from math import log
t = [2, 4, 6, 8, 12, 16]              # min
dP = [5.8, 9.9, 12.8, 14.9, 17.3, 18.6]  # kPa
dPmax = 19.8
for i in range(len(t)):
    k = log(dPmax / (dPmax - dP[i])) / t[i]
    print(t[i], round(k, 3))
\end{pythoncol}
\cote
\begin{consolecol}
2 0.173
4 0.173
6 0.173
8 0.175
12 0.172
16 0.175
\end{consolecol}
\end{codecote}
\item Vérifier la cohérence entre la constante $k$ et le temps de demi-réaction, sachant que $t_{1/2}=\ln2/k$ pour une loi
d'ordre 1. Calculer la vitesse volumique initiale de disparition du peroxyde d'hydrogène.
\item Sur une plaie, la décomposition est presque instantanée. Citer deux facteurs cinétiques qui l'expliquent.
\item Les flacons d'eau oxygénée sont opaques et doivent être conservés au frais. Justifier.
\end{questions}
\end{exoBac}

\begin{exoBac}{6}{Une lunette astronomique}
Un astronome amateur dispose d'une lunette dont l'objectif, de diamètre $D=\qty{70}{mm}$, a une distance focale
$f'_1=\qty{900}{mm}$, et de deux oculaires, de distances focales \qty{10}{mm} et \qty{25}{mm}. Il observe la planète Jupiter, dont le diamètre
apparent vaut 45 secondes d'arc.
\begin{donnees}
$\ang{1}=60'=3600''$. Pour les petits angles, $\tan\alpha\approx\alpha$ en radians. La diffraction limite la plus petite séparation angulaire
observable à $\theta\approx\dfrac{\lambda}{D}$, avec $\lambda=\qty{550}{nm}$. L'œil distingue deux points vus sous un angle d'au moins $1'$ environ.
\end{donnees}
\schema{bac-s4-lunette}
\begin{questions}
\item La lunette est afocale. Que signifie ce terme ? Où le foyer objet de l'oculaire doit-il être placé ? Calculer la
longueur de la lunette avec l'oculaire de \qty{10}{mm}.
\item Recopier le schéma et tracer la marche du faisceau à travers la lunette. Placer l'image intermédiaire $B_1$.
\item Définir le grossissement $G$ de la lunette et montrer que $G=\dfrac{f'_1}{f'_2}$.
\item Calculer le grossissement obtenu avec chaque oculaire.
\item Sous quel angle l'astronome voit-il Jupiter avec l'oculaire de \qty{10}{mm} ? Comparer au diamètre apparent de la Lune à l'œil nu,
voisin de \ang{0,5}.
\item \initiative Les bandes nuageuses de Jupiter ont une largeur apparente d'environ 4 secondes d'arc. Peuvent-elles être
distinguées avec cette lunette ? Un oculaire de \qty{4}{mm} permettrait-il d'en voir davantage de détails ?
\end{questions}
\end{exoBac}

\begin{exoBac}{5}{Le château d'eau}
Dans un village, l'eau potable est stockée dans un château d'eau ; la surface libre de l'eau, à l'air libre, se trouve à
$h=\qty{35}{m}$ au-dessus du robinet d'une maison.
\schema{bac-s4-chateau}
\begin{donnees}
$\rho=\qty{1000}{kg.m^{-3}}$ ; $g=\qty{9,81}{m.s^{-2}}$ ; $P_\mathrm{atm}=\qty{1,013e5}{Pa}$ ; diamètre intérieur du robinet : \qty{12}{mm}.
Relation de Bernoulli, pour un fluide parfait incompressible en écoulement permanent : $P+\frac12\,\rho\,v^2+\rho\,g\,z$ se conserve le
long d'une ligne de courant.
\end{donnees}
\begin{questions}
\item Le robinet étant fermé, l'eau est immobile. Calculer la pression de l'eau au niveau du robinet.
\item Pourquoi les châteaux d'eau sont-ils construits en hauteur ? Que se passe-t-il au dernier étage d'un immeuble de
\qty{40}{m} de hauteur alimenté par ce château d'eau ?
\item On ouvre le robinet. En appliquant la relation de Bernoulli entre la surface libre du château d'eau et la sortie du
robinet, calculer la vitesse de sortie de l'eau. Préciser les hypothèses.
\item En déduire le débit volumique prévu par ce modèle, en litres par minute.
\item \initiative Le débit mesuré n'est que de \qty{12}{L.min^{-1}}. Calculer la vitesse réelle de sortie de l'eau et commenter la validité du
modèle. Combien de temps faut-il pour remplir une baignoire de \qty{150}{L} ?
\end{questions}
\end{exoBac}
